\clearpage

\section{期末复习}

\subsection{复习提纲和建议}

\subsubsection{线性空间和线性映射（线代B）}

考试一般只出基础题，
要求包括：
知道如何证明一个空间是线性空间（封闭性+8条规则），
知道写出一个线性映射的表示矩阵（求出一组基在线性映射下的像，以列向量组的形式排成矩阵，按这个思路去写就不会搞错），
知道如何写出一个基变换的过渡矩阵（本质上是基变换的表示矩阵）.
坐标变换矩阵是过渡矩阵之逆.

\textbf{复习建议}：
（1）建议仔细学习书上的证明过程和例题，
即使看起来简单的证明也最好自己动手书写一下.
（2）本章的内容虽然在教科书里处于最后，
但事实上对四、五、六章的内容有统领性作用.
如果想要更深刻的理解线性代数的内容，
需要先把这一章吃透（详见讲义第四章）.

\subsubsection{知识串联：等价关系}

第四、五、六章中学习了矩阵的三种关系：
\textbf{相抵、相似和合同}.
它们都满足反身性、对称性和传递性.
我们称满足这三条性质的关系为\textbf{等价关系}.
对等价关系的详细叙述见讲义第二章2.1.2节.
须知谈论矩阵的某种等价关系下的性质，
即是对于矩阵进行\textbf{分类}，
又是在谈论矩阵在这种关系下的\textbf{不变量}，
还等同于在谈论\textbf{代表元/标准型}的性质.

以下从等价关系的不同视角，
解读一下矩阵的相抵、相似和合同关系.
要将这四个视角融为一体，
才算真正理解到位.
不过这要求是比较高的.
尤其是对于线代C的同学，
因为缺少第七、八章的知识，
所以对很多概念不了解是正常的.
从考试的角度来说，
尽力去掌握尽可能多的视角，
会让解题的道路更多，
并且会让书上的更多定理变得“显然”，
减少死记硬背的压力.



\medskip

\textbf{1. 相抵}

\textbf{从关系本身来看：}
$PAQ \sim A$，
其中$P,Q$是可逆矩阵.
也就是说可以通过左右各乘可逆矩阵进行转化的矩阵是相抵的.
或者说，
可以通过一系列初等行变换和列变换进行转化的矩阵是相抵的.

\textbf{从分类的角度来看：}
对于同一个线性映射，
在不同的基下可能表示成不同的矩阵，
这些矩阵相互之间是相抵的.
因此，
矩阵的相抵其实就是按照线性映射对矩阵（作为线性映射的表示矩阵）进行分类，
相互相抵的矩阵被归为同一类.

\textbf{从不变量的角度来看：}
两个矩阵相抵，当且仅当它们都是$m \times n$矩阵，且秩$r$相等.
这些变量构成了相抵关系的完全不变量.
它们的意义是给出了线性映射在同构意义下的本质信息（即与基的选取无关的性质）:
$n,m,r$分别是定义域、陪域和像空间的维数，
显然是与基的选取无关的.

\textbf{从相抵标准型的视角来看：}
任何矩阵都相抵于一个标准型$\mathrm{diag}(I_r, O)$.
标准型作为相抵等价类的代表，
一眼就能看出秩是多少，
而且形式简单，
便于处理.
在涉及到秩或者矩阵方程有解性（本质上也是在讨论矩阵的秩）时，
相抵的矩阵没有本质区别，
因此不妨拿出相抵标准型作为代表来处理问题.
具体题目详见讲义6.5.2节.

\bigskip

\textbf{2. 相似}

\textbf{从关系本身来看：}
$P^{-1} AP \sim A$，
其中$P$是可逆矩阵.
这个关系也可以从初等行列变换来理解.
如果说右乘矩阵$P$代表一系列初等列变换，
那么$P^{-1}$代表一系列对应的初等行变换之逆
（这种理解角度比较少用）.

\textbf{从分类的角度来看：}
对于同一个线性变换，
在不同的基下可能表示成不同的矩阵，
这些矩阵相互之间是相似的.
因此，
矩阵的相抵其实就是按照线性变换对矩阵（作为线性映射的表示矩阵）进行分类，
相互相抵的矩阵被归为同一类.
注意在写线性变换的表示矩阵时，
规定变换前后是用同一组基来展开的，
这造成了一个很大的限制.
因此相似的分类比相抵的分类要细致得多.

\textbf{从不变量的角度来看：}
相似不变量有：秩、行列式、
特征值、特征值的代数重数和几何重数、
特征多项式、迹.
它们都给出了矩阵背后的线性变换的性质（而不仅仅是表示矩阵的性质），
与基的选取无关.


\textbf{从标准型的视角来看：}
对于可对角化的矩阵，
其对角化的结果是这个线性变换的一种最简单的表示——
对角矩阵表示的是纯粹的伸缩变换（没有旋转或剪切）.
因此，
可对角化的矩阵本质上都可以看成简单的伸缩变换.
每个轴对应的伸缩系数就是特征值.

\bigskip

\textbf{3. 合同}

\textbf{从关系本身来看：}
$P^T AP \sim A$，
其中$P$是可逆矩阵.
从初等行列变换的角度来理解：
成对的初等行列变换就是合同变换.

\textbf{从分类的角度来看：}
对于同一个广义内积结构，
在不同的基下可能表示成不同的度规矩阵（实对称矩阵），
这些矩阵相互之间是合同的.
因此，
矩阵的相抵其实就是按照广义内积的性质对矩阵（作为度规矩阵）进行分类，
相互合同的矩阵被归为同一类.
其中，实内积结构（满足内积的正定性，即非零向量的长度平方必然为正数）
的度规矩阵是单独一类，
称为正定矩阵.

\textbf{从不变量的角度来看：}
秩和正惯性指数构成了合同关系的完全不变量.
因为着秩和正惯性指数是广义内积结构的本质的性质，
与基的选取无关.
它反映了一个广义内积会把空间里“多少维”的向量指定为类空、类光或者类时的.
正定矩阵所反映的实内积结构会将整个空间指定为类空区域.

\textbf{从标准型的视角来看：}
任何实对称矩阵总能合同变换为合同标准型（对角矩阵），
以及唯一的规范型（对角元只有$1,-1,0$）.
规范型的非零元个数就是秩，
$1$的个数就是正惯性指数.
正定矩阵即是规范型为$I_n$的矩阵.


\subsubsection{矩阵的运算}

本章复习的总体建议：
尽量去从子空间和线性映射的角度来理解矩阵的各个性质和运算.
如果做不到，
那么需要仔细阅读教材上叙述的一些“小性质”.
比如在矩阵的各种运算（加法、数乘、矩阵乘法、逆、转置）下，
秩怎么变，行列式怎么变？
可以记住一些教材课后习题以及讲义习题中，一些秩的等式和不等式，
作为二级结论备用.

\medskip

\textbf{1. 理解左（右）乘可逆矩阵等同于做一系列初等行（列）变换}.
这个知识点是贯穿下半学期的.
建议仔细了解这一结论的来龙去脉，
学习书上的证明过程，
并以此来深刻理解相抵关系.

\textbf{计算题}：
求解矩阵方程（如$AX=B, AXB = C$）.
较快的方法是把未知矩阵$X$理解为一系列初等行列变换，
用初等变换法去解.
同样的方法可以用来求逆矩阵.

\medskip

\textbf{2. 可逆矩阵相关}.
求矩阵的逆
有如下几种办法：

(1)~初等变换法，如上文所述.
对于具体给定的矩阵，
一般来说是最快的方法.

(2)~伴随矩阵法：利用$AA^* = |A|I$.
一般来说不建议用来求矩阵的逆.
但是题目有可能直接涉及伴随矩阵的性质.
伴随矩阵相关的性质和处理办法，
结合讲义5.4.4节复习.

(3)~凑.讲义5.4.3与矩阵求逆有关的例题.
当矩阵难以直接求逆，
但矩阵本身很有规律时，
可以通过矩阵分解的方式表达出你所看到的规律，
然后根据分解式用待定系数法去凑.

(4)~正交矩阵.
如果矩阵是个正交矩阵，
那么将其转置就是逆矩阵.
注意捡漏.

如果是判断一个矩阵是否可逆，
思路要灵活.
最基本的是知道矩阵可逆$\Leftrightarrow$矩阵满秩$\Leftrightarrow$行列式非零$\Leftrightarrow$相抵于单位矩阵.
但是注意讨论相似关系和合同关系时也可以“向下”得到矩阵的秩.
比如矩阵的秩同时也是非零特征值的个数，
因此如果矩阵有$n$个非零特征值，
也可以判断矩阵可逆.
另外，
秩也是合同不变量.
特别地，正定（负定）矩阵一定是可逆的.
这些地方考验对概念的理解程度，
光靠背诵难以全数掌握.

\medskip

\textbf{3. 分块矩阵}.
分块矩阵的运算，以及分块初等行列变换，
是非常重要的理论和计算工具.
需要仔细阅读教材上的相关叙述，
尤其是做初等行列变换时是否需要乘可逆矩阵，
是左乘还是右乘，
一定要搞清楚.
在解题时，
分块矩阵对于研究矩阵的秩有大用，
也是线代B和线代C非基础题型的常客.
结合讲义5.1.2和讲义5.4.1的习题来复习.

\subsubsection{矩阵的相似}

总体复习建议：
吃透矩阵对角化的计算过程，
理解每一步的原因是什么？
如果矩阵不可对角化，
还强行尝试去计算对角化，
会在哪一步遇到阻碍？
想清楚这些问题，
等于是搞明白了本章中矩阵可对角化的条件（不包括实对称矩阵的特殊性质），
不需要专门背诵.

有关矩阵可对角化条件的定理，
不建议学习书上的证明过程
（尤其是实对称矩阵相关的，直接背诵“实对称矩阵一定可以在实数范围内正交对角化”即可）.
如果你熟悉第七、八章的内容，
可以思考是否有更加简洁明了的证明方式.

\medskip

\textbf{1. 相似不变量}
以特征多项式为核心.
特征多项式可以衍生出几乎所有（教材提及的）其他相似不变量.
特征多项式是线性代数范围内“最强大的”相似不变量.

但是注意，
本章对于相似关系的研究是很浅的.
对于不可对角化的矩阵，
我们几乎没有手段去研究，
连判断任意两个矩阵是否相似都很困难.

\medskip

\textbf{2. 矩阵可对角化的条件}.
讲义这部分写得比较精简，
我这里不额外列复习提纲了.
从解题的角度来说，
注意$\lambda$的几何重数就是$\dim \operatorname{Ker}(\lambda I - A) = n - \mathrm{rank}(\lambda I - A)$.
因此，“几何重数之和为$n$”的条件实际上给出了一个秩等式.
而秩等式的证明用上一章讲的分块矩阵法/线性映射法可以套路化地解决.

\subsubsection{矩阵的合同}

总体复习建议：
由于秩和正惯性指数给出了合同等价关系的完全不变量，
因此合同关系在线性代数的范围里是被研究透彻的了.
因此比较考验对等价关系的理解.
以正定矩阵为例，
能否将这个概念从二次型，合同规范型，正惯性指数的角度“无痛转化”进行叙述？

本章教材中的证明多数都不复杂，
看一看会有启发.
尤其是数学归纳法的思路很有学习价值（例如证明正定等价于顺序主子式全都是正数.）


\medskip

\textbf{1. ~合同与相似的关联——正交对角化}.
对于实对称矩阵$A$，
总可以找到正交矩阵$P$，
使得$P^{-1}AP$为对角矩阵$\mathrm{diag}(\lambda_1 \cdots , \lambda_n)$.
注意正交矩阵满足$P^{-1} = P^T$，
因此得到的这个相似标准型其实也是一个合同标准型.
因此，
从实对称矩阵的相似标准型（即从$n$个特征值中）也可以"读"出合同的性质.
比如，
实对称矩阵$A$的正惯性指数就是正特征值的个数；
实对称矩阵$A$正定当且仅当所有特征值都是正数.
这里就建立了合同与相似之间的桥梁.

\textbf{2.}
判断一个矩阵是否正定，
除了从定义本身出发，
也可以证明它的规范型是单位矩阵.
在计算题中，
常用的方法是去计算顺序主子式.






\subsection{经典解题思路补充}

\subsubsection{矩阵运算}

\begin{example}
    设矩阵：
    \begin{eq1}
        A = 
        \begin{pmatrix}
            1 & \bm{\alpha} & \bm{\beta} \\
            0 & 1 & \bm{\alpha} \\
            0 & 0 & 1
        \end{pmatrix}
    \end{eq1}
    计算$A^n$.
\end{example}

\bigskip

\begin{example}
    设$A$是已知的$n$阶方阵，
    且$A^3 = 2I$.
    $B = A^2 - 2A + 2I$.
    证明：$B$是满秩矩阵，
    并求$B^{-1}$.
\end{example}

\begin{example}
    设 $A$ 为 $m \times n$ 矩阵，
    $B$ 为 $n \times m$ 矩阵，
    使得 $I_m + AB$ 可逆，求证：
    $I_n + BA$也可逆.
\end{example}

\bigskip

\begin{example}
    设$A,B$是$n$阶方阵，
    $\mathrm{rank}(I-AB) + \mathrm{rank}(I+BA) = n$.
    求证：$\mathrm{rank}(A) = n$.
\end{example}

\bigskip

\begin{example}
    求下列 $n$ 阶矩阵的逆阵，
    其中 $a_i \neq 0 \ (1 \leq i \leq n)$：
\begin{eq1}
A =
\begin{pmatrix}
1 + a_1 & 1 & 1 & \cdots & 1 \\
1 & 1 + a_2 & 1 & \cdots & 1 \\
1 & 1 & 1 + a_3 & \cdots & 1 \\
\vdots & \vdots & \vdots & \ddots & \vdots \\
1 & 1 & 1 & \cdots & 1 + a_n
\end{pmatrix}.
\end{eq1}

\end{example}


\bigskip
\bigskip

\begin{example}
    证明：
    与$n$阶Jordan块
    \begin{eq1}
        J = 
        \begin{pmatrix}
            \lambda & 1 & ~ & ~  \\
            ~ & \lambda & \ddots ~ \\
            ~ & ~ & \ddots & 1 \\
            ~ & ~ & ~ & \lambda
        \end{pmatrix}
    \end{eq1}
    可交换的矩阵必然是$J$的多项式.
\end{example}

\bigskip

\begin{example}
    设 $n$ 阶方阵 $A$ 的每一行元素之和等于常数 $c$，求证：

    (1)~对任意的正整数 $k$，$A^k$ 的每一行元素之和等于 $c^k$.

    (2)~若 $A$ 为可逆阵，则 $c \neq 0$ 且 $A^{-1}$ 的每一行元素之和等于 $c^{-1}$.
\end{example}

\bigskip

\begin{example}
    设 $A$ 为 $m$ 阶矩阵，
    $B$ 为 $n$ 阶矩阵，
    求分块对角阵 $C$ 的伴随矩阵：
    \begin{eq1}
    C =
    \begin{pmatrix}
    A & O \\
    O & B
    \end{pmatrix}.
    \end{eq1}
\end{example}

\subsubsection{矩阵的相抵与相似}

\begin{example}
    讲义习题6.4证明了满秩分解定理：
    设$A_{m \times n}$的秩为$r$.
    证明：存在秩为$r$的矩阵$B_{m \times r}$和秩为$r$的矩阵$C_{r \times n}$，
    使得$A = BC$.
    那么，
    对于任给的矩阵$A$，
    如何计算它的满秩分解，
    即如何求出$B$和$C$？
    以下面的矩阵$A$为例展示具体的计算步骤：
    \begin{eq1}
        A = 
    \begin{pmatrix}
        1 & -1 & -3 & 1 \\
        1 & -1 & 2 & -1 \\
        4 & -4 & 3 & -2 
    \end{pmatrix} 
    \end{eq1}
\end{example}

\bigskip

\begin{example}
    证明：

    (1)~对任意矩阵$A$，矩阵方程$AXA = A$都有解.

    (2)~如果矩阵方程$AY=C$和$ZB=C$有解，
    那么$AXB = C$有解.
\end{example}

\bigskip

\begin{example}
    求出所有仅与自身相似的矩阵.
\end{example}

\bigskip

\begin{example}
    设$A$是$n$阶方阵，
    且存在$n$维列向量$\bm{\alpha}$，
    使向量组$\bm{\alpha} , A \bm{\alpha} , A^2 \bm{\alpha} ,\cdots, A^{n-1}\bm{\alpha}$线性无关.
    证明：$A$的任何特征子空间必定只有$1$维.
\end{example}

\bigskip

\begin{example}
    设$A$是$n$阶方阵，
    $\lambda$是它的一个特征值，
    $\bm{\alpha}$是属于$\lambda$的一个特征向量.
    证明：存在数$\mu$，
    使得$\mu$是$A^*$的一个特征值，
    $\bm{\alpha}$的$A^*$的属于特征值$\mu$的一个特征向量.
\end{example}

\bigskip

\begin{example}
    设$A,B$是$n$阶矩阵，
    满足$AB=BA$.

    (1)~若$A,B$都是复矩阵.
    证明：$A,B$至少有一个公共的特征向量.

    (2)~上述结论对于其他数域也都成立吗？
\end{example}

\bigskip

\begin{example}
    ~

    (1)~思考：对矩阵$A$做相似变换$A \mapsto P^{-1}AP$，
    相当于对$A$做什么样的初等行列变换？

    (2)~设$n$阶方阵$A,B$满足$\mathrm{rank}(ABA) = \mathrm{rank}(B)$.
    求证：$AB$与$BA$相似.
\end{example}

\bigskip


\subsubsection{矩阵的合同}

\begin{example}
    设$f = \bm{X}^T A \bm{X}$是非退化的二次型，
    其中$A$是对称矩阵.
    证明：$f$可用正交变换化为规范型，当且仅当$A$是正交矩阵.
\end{example}

\begin{example}
    设$A$是$n$阶正定实对称矩阵，
    $A$的非对角元均小于零.
    证明：$A$是可逆矩阵，
    且$A^{-1}$的每个元素均大于零.
\end{example}

\begin{example}
    设$A$是$n$阶可逆实对称矩阵.
    证明：$A$为正定矩阵当且仅当对所有$n$阶正定矩阵$B$，
    有$\mathrm{tr}(AB) > 0$.
\end{example}

\begin{example}
    秩等于$r$的对称矩阵等于$r$个秩为$1$的对称矩阵之和.
\end{example}

\begin{example}
    设 $A$ 为 $n$ 阶实对称矩阵，求证：

    (1)~$A$ 是正定阵的充要条件是存在 $n$ 阶非异实矩阵 $C$，使得 $A = C^T C$.
    
    (2)~$A$ 是半正定阵的充要条件是存在 $n$ 阶实矩阵 $C$，使得 $A = C^T C$.
\end{example}

\begin{example}
    证明：一个秩大于 $1$ 的实二次型可以分解为两个实系数一次多项式之积的充要条件是它的秩等于 $2$，且符号差等于零.
\end{example}

\bigskip

\subsubsection{线性空间与线性映射}

\begin{example}
    设$V$是所有$n$阶实方阵构成的$n^2$维向量空间，
    $U$和$W$分别为所有$n$阶实对称矩阵和反对称矩阵构成的子空间.
    证明：$V = U \oplus W$.
\end{example}

\bigskip

\begin{example}
    设$n$维向量空间$V$上的线性变换$\mathcal{P}$满足$\mathcal{P}^2 = \mathcal{P}$.
    证明：$V  =  \operatorname{Im}\mathcal{P} \oplus \operatorname{Ker} \mathcal{P}$.
\end{example}

\bigskip

\begin{example}
    设 $V$ 是次数不超过 $n$ 的实系数多项式全体组成的线性空间，求从基
$\{1, x, x^2, \cdots, x^n\}$
到基$\{1, x-a, (x-a)^2, \cdots, (x-a)^n\}$
的过渡矩阵，并以此证明多项式的 Taylor 公式：
\begin{eq1}
f(x) = f(a) + \frac{f'(a)}{1!}(x-a) + \frac{f''(a)}{2!}(x-a)^2 + \cdots + \frac{f^{(n)}(a)}{n!}(x-a)^n,
\end{eq1}
其中 $f^{(n)}(x)$ 表示 $f(x)$ 的 $n$ 次导数。
\end{example}

\subsection{其他拓展题目}


\begin{example}
    一个行列式为$1$的$n$阶方阵能否写成若干个行列式为$1$的初等矩阵之积？
\end{example}

\bigskip

\begin{example}
    我们知道，
    若$n$阶矩阵$A$可对角化，
    则可以利用对角化方便地计算$A^n$.
    但如果$A$不可对角化呢？
    以下面的矩阵为例：
    \begin{eq1}
        A = 
        \begin{pmatrix}
            4 & 2 & 2 \\
            0 & 4 & 0 \\
            0 & -2 & -2 
        \end{pmatrix}
    \end{eq1}

    (1)~计算它的全部特征值和特征向量，
    并说明它为什么不可对角化.

    (2)~求可逆矩阵$P$，
    使得$P^{-1}AP = J$.
    其中$J$是$A$的Jordan标准型：
    \begin{eq1}
        J = 
        \begin{pmatrix}
            -2 & ~ & ~ \\
            ~ & 4 & 1 \\
            ~ & ~ & 4
        \end{pmatrix}
    \end{eq1}

    (3)~计算$A^{20}$.
\end{example}

\bigskip

\begin{example}
    设$f: M_{n}(\mathbb{K}) \to \mathbb{K}$是线性映射，
    且满足交换性和正规性：

    (i)~交换性：$f(AB) = f(BA)$.

    (ii)~正规性：$f(I_n) = n$.

    证明：$f$就是迹.
    即$f(A) = \mathrm{tr}(A)$.
\end{example}

\bigskip

\begin{example}~

    (1)~设$A$是数域$\mathbb{K}$上的$n$阶方阵，
    且它有$n$个$\mathbb{K}$中的特征值.
    证明：存在数域$\mathbb{K}$上的可逆矩阵$P$，
    使得$P^{-1}AP$是上三角矩阵.

    (2)~
    证明：任意$n$阶复矩阵相似于上三角矩阵.
\end{example}